% BOOK-FILE-ID: f-947cf9e311ee
\subsubsection{整数から作る近似分数}\label{節：整数から作る近似分数}
\sectionref*{節：ニュートン法}では,$a_1=2$から$\sqrt2$を近似した.
ここでは同じ漸化式$a_{n+1}=\frac12\left(a_n+\frac{2}{a_n}\right)$で,初項を$a_1=1$に変えてみる.
\[
  a_2=\frac32,\qquad a_3=\frac{17}{12},\qquad a_4=\frac{577}{408}
\]
となる.初項$1,2$のどちらも$\frac{3}{2}$へ移るため,第2項以降は本編と一致する.
分子と分母を別々に調べ,連立型・一次分数型・\bookterm{newton}{ニュートン法}をつなげる.微積分は使わない.

\paragraphbf{分子と分母の関係}
$3^2-2\cdot2^2=17^2-2\cdot12^2=577^2-2\cdot408^2=1$である.
次の\bookterm{sequence}{数列}を作ろう.
\[
  x_k+y_k\sqrt2=(3+2\sqrt2)^k\quad(k\ge1)
\]
$3+2\sqrt2$倍して係数を比べると,
\[
  x_{k+1}=3x_k+4y_k,\qquad y_{k+1}=2x_k+3y_k
\]
となる.$x_1=3$, $y_1=2$なので,両方とも正の整数である.
$\sqrt2$の符号を変えた式との積から
\[
  x_k^2-2y_k^2=(3+2\sqrt2)^k(3-2\sqrt2)^k=1
\]
が保たれる.整数の式を保つ漸化式は,
\bookproblemref{practice-bridge}{9}{発展入門問9}の考え方とも共通する.

\paragraphbf{比から一次分数型へ}
$r_k=\frac{x_k}{y_k}$とすると,連立漸化式から
\[
  r_{k+1}=\frac{3r_k+4}{2r_k+3}
\]
である.分母は正であり,\bookterm{fixed}{不動点}を求めると$\pm\sqrt2$になる.
一次分数型で用いた差の比は,
\[
  \frac{r_{k+1}-\sqrt2}{r_{k+1}+\sqrt2}
  =(3-2\sqrt2)^2
  \frac{r_k-\sqrt2}{r_k+\sqrt2}
\]
となる.整数の組とその比は,同じ漸化式の2通りの表し方である.

\paragraphbf[例：次の近似分数]{例：$\frac{17}{12}$の次の近似分数}
本編で使った漸化式に$\frac{17}{12}$を代入すると,
\[
  \frac12\left(\frac{17}{12}+\frac{24}{17}\right)=\frac{577}{408}
\]
となる.一方,整数の連立漸化式を1回用いた$r_3$は$\frac{99}{70}$である.
ニュートン法は,この列で次の分数へ一つ進んでいるわけではない.

\par\medskip\noindent\textbf{別解：ニュートン法と添字の倍増}\par\nobreak\smallskip
$(x_k+y_k\sqrt2)^2$の展開から,
\[
  x_{2k}=x_k^2+2y_k^2,\qquad y_{2k}=2x_ky_k
\]
なので,
\[
  r_{2k}=\frac12\left(r_k+\frac2{r_k}\right).
\]
この\bookterm{newton}{ニュートン法}型の式から
\[
  a_n=r_{2^{n-2}}\qquad(n\ge2)
\]
となる.$\frac{3}{2}$, $\frac{17}{12}$, $\frac{577}{408}$は$r_1$, $r_2$, $r_4$であり,連続する三項ではない.
\BookFigPellDoubling
$x_k^2-2y_k^2=1$を因数分解すれば,
\[
  r_k-\sqrt2=\frac1{y_k^2(r_k+\sqrt2)}>0.
\]
$y_{k+1}>3y_k$なので$y_k\to\infty$であり,$r_k\to\sqrt2$となる.
添字の倍増からも,ニュートン法の速さを読める.

\paragraphbf{ペル方程式}
平方数でない正の整数$d$に対し,$x^2-dy^2=1$の整数解を求める問題を\textbf{\bookterm{pell}{ペル方程式}}という.\footnote{ここでは$d=2$の具体例を扱い,一般の解の存在や全解の分類には進まない.}

\paragraphbf{連分数にも現れる分数}
$\sqrt2=1+\frac1{1+\sqrt2}$から,有限段階の近似を
\[
  C_0=1,\qquad C_{j+1}=1+\frac1{1+C_j}
  \quad(j\ge0)
\]
で定める.$C_1=\frac{3}{2}$, $C_2=\frac{7}{5}$, $C_3=\frac{17}{12}$である.
$C_j=\frac{P_j}{Q_j}$とすると,$P_0=Q_0=1$から
\[
  P_{j+1}=P_j+2Q_j,\qquad Q_{j+1}=P_j+Q_j
\]
が得られ,$P_j+Q_j\sqrt2=(1+\sqrt2)^{j+1}$となる.
$(1+\sqrt2)^2=3+2\sqrt2$だから,
\[
  C_{2k-1}=r_k\qquad(k\ge1)
\]
である.この有限段階の値を\textbf{\bookterm{continued}{連分数}の近似分数}と呼ぶ.
偶数番号では$P_j^2-2Q_j^2=-1$,奇数番号では$1$となり,
同じ誤差の因数分解で$C_j\to\sqrt2$も確認できる.

\begin{bookpoint}
  \textbf{POINT}\quad 三つの解法のつながり\par
  分子・分母を調べれば連立型,比を調べれば一次分数型,
  添字を倍に進めれば平方根の\bookterm{newton}{ニュートン法}になる.
  置く量と漸化式を用いる回数を変えると,解法どうしが結び付く.
\end{bookpoint}
保存される式と誤差を\bookproblemref{practice-connections}{6}{接続演習問6}で確かめよう.
\bookterm{pell}{ペル方程式}と\bookterm{continued}{連分数}の補足資料は,\chapterref*{章：参考文献}のKeith Conradの二つの解説にある.
